\chapter{Стационарный режим теплопроводности}

\section{Определение стационарного режима}

Стационарным режимом называется распределение температуры, не зависящее
от времени:
\[
 u(x,t)=U(x),\qquad u_t(x,t)=0.
\]
Температура в каждой точке стержня остаётся постоянной, однако это
не означает отсутствия тепловых потоков. Потоки, теплоотдача и источники
тепла взаимно компенсируются, так что накопления тепла не происходит.

Для задачи
\[
\begin{cases}
 u_t=a^2u_{xx}-bu+f(x,t),&0<x<l,\\
 u(x,0)=\phi(x),\\
 \alpha u(0,t)-\beta u_x(0,t)=\mu(t),\\
 \gamma u(l,t)-\delta u_x(l,t)=\nu(t)
\end{cases}
\]
точный стационарный режим возможен лишь при независимых от времени
правых частях:
\[
 f(x,t)=f_s(x),\qquad\mu(t)=\mu_s,\qquad\nu(t)=\nu_s.
\]
Действительно, подстановка $u=U(x)$ делает левые части уравнения
и граничных условий независимыми от времени.
Профиль $U$ определяется краевой задачей
\[
\boxed{\begin{cases}
 a^2U''(x)-bU(x)+f_s(x)=0,&0<x<l,\\
 \alpha U(0)-\beta U'(0)=\mu_s,\\
 \gamma U(l)-\delta U'(l)=\nu_s.
\end{cases}}
\]
Начальное условие в эту краевую задачу не входит.
Если $\phi=U$, то решение стационарно с начального момента;
если $\phi\ne U$, возникает переходный процесс.

Следует различать существование стационарного решения и установление
стационарного режима:
\[
 \lim_{t\to\infty}u(\cdot,t)=U.
\]
Для второго утверждения необходима устойчивость стационарного решения.
Если данные лишь стремятся к $f_s,\mu_s,\nu_s$, то предельный профиль
определяется той же краевой задачей; сходимость требует дополнительных
условий устойчивости и регулярности данных.

\section{Влияние неоднородности уравнения}

Функция $f_s(x)$ описывает распределённое тепловое воздействие:
внутренние источники и, при принятом ранее переобозначении,
вклад температуры окружающей среды. Она непосредственно определяет
форму стационарного профиля:
\[
 U''(x)=\frac{b}{a^2}U(x)-\frac{f_s(x)}{a^2}.
\]
При $b=0$ имеем
\[
 U''(x)=-\frac{f_s(x)}{a^2}.
\]
Поэтому при $f_s=0$ профиль линейный, а при постоянном ненулевом
источнике --- квадратичный. При $b>0$ даже при $f_s=0$ профиль
в общем случае является комбинацией гиперболических функций.

Для условий первого рода при $x=0$ и второго рода при $x=l$ положим
\[
 U(0)=T_0,\qquad U'(l)=q_0.
\]
Здесь $q_0$ --- заданная производная температуры, а не непосредственно
тепловой поток. Поток в положительном направлении оси равен
$-kS q_0$, где $k$ --- теплопроводность и $S$ --- площадь сечения.
При $b=0$ интегрирование даёт
\[
\boxed{U(x)=T_0+q_0x+
 \frac{1}{a^2}\int_0^l\min(x,\xi)f_s(\xi)\,d\xi.}
\]
В частности, если $f_s(x)=F_0$, то
\[
 U(x)=T_0+q_0x+\frac{F_0}{a^2}
 \left(lx-\frac{x^2}{2}\right).
\]
Таким образом, изменение источника изменяет стационарный профиль
даже при неизменных данных на границах.

\section{Влияние начальных условий}

Пусть стационарный профиль $U$ существует. Выполним замену
\[
 u(x,t)=U(x)+v(x,t).
\]
При постоянных во времени внешних данных функция $v$ удовлетворяет
\[
 v_t=a^2v_{xx}-bv,\qquad v(x,0)=\phi(x)-U(x)
\]
с однородными граничными условиями.
Пусть $X_n$ --- собственные функции задачи
\[
 -X_n''=\lambda_nX_n,\qquad
 \alpha X_n(0)-\beta X_n'(0)=0,\qquad
 \gamma X_n(l)-\delta X_n'(l)=0.
\]
Тогда
\[
 v(x,t)=\sum_n c_n e^{-(a^2\lambda_n+b)t}X_n(x),
 \qquad
 c_n=\frac{\int_0^l(\phi-U)X_n\,dx}{\int_0^lX_n^2\,dx}.
\]
Если для всех мод выполняется
\[
 a^2\lambda_n+b>0,
\]
переходная часть затухает. В этом случае стационарный профиль единственен
и не зависит от начальных условий; начальная температура определяет
амплитуды переходных мод и время достижения заданной близости к равновесию.

Для рассматриваемых условий первого--второго рода
\[
 \lambda_n=\left(\frac{(n+\tfrac12)\pi}{l}\right)^2,
 \qquad n=0,1,2,\ldots.
\]
При $b\ge0$ все моды затухают. В норме $L^2(0,l)$
\[
 \|u(\cdot,t)-U\|_{L^2}
 \le e^{-\sigma_0t}\|\phi-U\|_{L^2},\qquad
 \sigma_0=b+\frac{a^2\pi^2}{4l^2}>0.
\]

Утверждение о независимости предельного режима от начальных данных
не является универсальным. Если оба конца теплоизолированы,
$b=0$ и $f_s=0$, то средняя температура сохраняется и
\[
 u(x,t)\longrightarrow\frac1l\int_0^l\phi(\xi)\,d\xi.
\]
Здесь существует семейство постоянных стационарных решений, а начальные
условия выбирают одно из них. Если имеется незатухающая или растущая мода,
общая сходимость к единственному стационарному режиму отсутствует.
В частности, для записи с минусом перед производной на обоих концах
не следует предполагать положительность спектра без проверки.

\section{Влияние граничных неоднородностей}

Постоянные значения $\mu_s$ и $\nu_s$ непосредственно входят в
стационарную краевую задачу, поэтому определяют уровень температуры
и потоки на концах стержня. Даже если $f_s=0$, неоднородные граничные
условия могут поддерживать ненулевой стационарный профиль.

При выбранных коэффициентах
\[
 \alpha>0,\quad\beta=0,\qquad\gamma=0,\quad\delta>0
\]
получаем
\[
 T_0=\frac{\mu_s}{\alpha},\qquad
 q_0=-\frac{\nu_s}{\delta}.
\]
При $f_s=0$ и $b=0$
\[
\boxed{U(x)=\frac{\mu_s}{\alpha}-\frac{\nu_s}{\delta}x.}
\]
При $f_s=0$ и $b>0$, обозначая $\kappa=\sqrt{b}/a$, имеем
\[
\boxed{U(x)=T_0\cosh(\kappa x)+
 \frac{q_0/\kappa-T_0\sinh(\kappa l)}{\cosh(\kappa l)}
 \sinh(\kappa x).}
\]
Эта формула показывает совместное влияние граничных данных
и теплоотдачи через боковую поверхность.

При переменных во времени $\mu(t)$ и $\nu(t)$ точный стационарный
профиль невозможен. Если они стремятся к постоянным значениям
и система устойчива, возможно установление предельного стационарного
режима. Периодическое граничное воздействие, напротив, может приводить
к установившемуся периодическому режиму, который не является стационарным.

\section{Баланс и существование режима}

Интегрирование стационарного уравнения даёт необходимый баланс:
\[
 a^2\bigl(U'(l)-U'(0)\bigr)-b\int_0^lU(x)\,dx+
 \int_0^lf_s(x)\,dx=0.
\]
Это равенство выражает отсутствие накопления тепла:
поступление через источники и границы компенсируется потерями.
Оно не заменяет полного решения краевой задачи.

При теплоизоляции обоих концов и $b=0$ необходимо
\[
 \int_0^lf_s(x)\,dx=0.
\]
Если это условие нарушено, средняя температура изменяется со временем,
и стационарного режима нет. При выполнении условия профиль определяется
с точностью до аддитивной константы, выбираемой начальной средней
температурой.

Для условий первого рода слева и второго рода справа при $b\ge0$
стационарная задача имеет единственное решение для достаточно регулярных
$f_s$ и любых постоянных граничных данных. Этот профиль определяется
источниками, граничными условиями и коэффициентами модели;
начальные данные влияют только на переходный процесс.
